\scanpage{75}
% slanted
\reversemarginpar\margin{\shortstack[l]{\unsure{Attach|attach}\\SGA 7 \(\overline{\underline{\mathrm{VI}}}\) ?\\Appendix \(\overline{\underline{\mathrm{II}}}\)}}\normalmarginpar
\noindent\hfill Bures sur Yvette, 21 October 1968\\

\noindent Dear Grothendieck,\\

In fact, it is not necessary to carry out elaborate gluings of cotangent complexes to prove that

\noindent \uline{proposition} : \uline{Let} \(X\) \uline{be a compact complex analytic space}. \uline{For} \(X\) \uline{to be a scheme}, \uline{it is necessary and sufficient that} \(X_{\mathrm{red}}\) \uline{be one}

One reduces to \deleted{assuming} showing that if \(X_0\) is defined by an Ideal of square zero, say \(\mathscr{N}\), of \(X\), then
\[
 (X_0\text{ is a scheme})\ \Longrightarrow\ (X\text{ is a scheme})
\]

Let \(E^{\mathrm{an}}(X_0,\mathscr{N})\) and \(E^{\mathrm{alg}}(X_0,\mathscr{N})\) be the analytic and algebraic groups of infinitesimal extensions of \(X_0\) by \(\mathscr{N}\). Finally let \(\unsure{K|\mathbb{K}}\) be the cotangent complex of \(X_0\), \uline{given locally, in the derived category}. it is obtained by embedding \uline{locally} (Zariski, or étale) \(X_0\) in a smooth \unsure{scheme} \ldots\ and allows one to define
\begingroup
\makeatletter\tagsleft@true\makeatother
\[
 \mathscr{H}om(K,\mathscr{N})\qquad\text{and}\qquad
 \mathscr{E}xt^{1}\makebox[0pt][l]{\hspace{-0.20em}\raisebox{-1.8ex}[0pt][0pt]{\parbox[c]{1.5mm}{\figure{\unsure{%
 \begin{tikzpicture}[x=1mm,y=1mm,baseline=0pt]
 \fill[gray!75] (-0.28,-0.02) .. controls (-0.28,0.12) and (-0.11,0.27) .. (0.05,0.25)
 .. controls (0.19,0.23) and (0.29,0.06) .. (0.22,-0.03) -- cycle;
 \fill[gray!30] (0.62,0.08) ellipse[x radius=0.08,y radius=0.17,rotate=32];
 \end{tikzpicture}%
 }}}}}(K,\mathscr{N})\ .\tag*{(1)}
\]
\endgroup

Consider the diagram :

\margin{\vspace*{1.4em}\figure{\unsure{%
\begin{tikzpicture}[x=1mm,y=1mm]
\fill (-0.22,-0.07)
 .. controls (-0.27,0.06) and (-0.16,0.16) .. (-0.08,0.17)
 .. controls (-0.04,0.30) and (0.06,0.25) .. (0.11,0.14)
 .. controls (0.25,0.10) and (0.27,-0.03) .. (0.15,-0.11)
 .. controls (0.04,-0.15) and (-0.15,-0.16) .. (-0.22,-0.07) -- cycle;
\fill (0.37,0.72)
 .. controls (0.34,0.85) and (0.44,0.95) .. (0.50,0.98)
 .. controls (0.54,1.10) and (0.65,1.11) .. (0.71,1.00)
 .. controls (0.83,0.94) and (0.77,0.85) .. (0.68,0.81)
 .. controls (0.73,0.69) and (0.64,0.62) .. (0.53,0.65)
 .. controls (0.44,0.63) and (0.38,0.65) .. (0.37,0.72) -- cycle;
\fill[gray!65] (0.20,-0.19) .. controls (0.22,-0.13) and (0.31,-0.23) .. (0.34,-0.34)
 .. controls (0.30,-0.46) and (0.18,-0.36) .. (0.20,-0.19) -- cycle;
\fill[gray!50] (0.47,-0.50) ellipse[x radius=0.06,y radius=0.08];
\end{tikzpicture}%
}}}

\diagram{\footnotesize
\begin{tikzcd}[ampersand replacement=\&,column sep=0.85em,row sep=1.6em]
 (2) \& 0 \arrow[r] \&
 H^{1^{\mathrm{alg}}}(X_0,\mathscr{H}om(K,\mathscr{N})) \arrow[r] \arrow[d] \&
 E^{\mathrm{alg}}(X_0,\mathscr{N}) \arrow[r,"\alpha"] \arrow[d] \&
 H^{0}(X_0,\mathscr{E}xt^{1}(K,\mathscr{N})) \arrow[d] \\
 \& 0 \arrow[r] \&
 H^{1^{\mathrm{an}}}(X_0,\mathscr{H}om(K,\mathscr{N})) \arrow[r] \&
 E^{\mathrm{an}}(X_0,\mathscr{N}) \arrow[r,"\beta"] \&
 H^{0^{\mathrm{an}}}(X_0,\mathscr{E}xt^{1}(K,\mathscr{N}))
\end{tikzcd}}


In view of the (local) meanings of (1), the rows are exact ; by GAGA, the outer vertical arrows are isomorphisms. It remains to \deleted{compare the images of \mbox{\(\alpha\)} and \mbox{\(\beta\)}} show that \(\operatorname{Im}(\beta)\subset\operatorname{Im}(\alpha)\).
% archive sheet number: 74
\clearpage
